%==============================================================================
% Discussion Intro (draft scratchpad)
%==============================================================================
%
% Use this file to draft/iterate on the opening of the Discussion chapter.
% It is NOT included anywhere by default.
%
% If you want to include it in the thesis build later, add one of:
%   %==============================================================================
% Discussion Intro (draft scratchpad)
%==============================================================================
%
% Use this file to draft/iterate on the opening of the Discussion chapter.
% It is NOT included anywhere by default.
%
% If you want to include it in the thesis build later, add one of:
%   %==============================================================================
% Discussion Intro (draft scratchpad)
%==============================================================================
%
% Use this file to draft/iterate on the opening of the Discussion chapter.
% It is NOT included anywhere by default.
%
% If you want to include it in the thesis build later, add one of:
%   %==============================================================================
% Discussion Intro (draft scratchpad)
%==============================================================================
%
% Use this file to draft/iterate on the opening of the Discussion chapter.
% It is NOT included anywhere by default.
%
% If you want to include it in the thesis build later, add one of:
%   \input{Chapters/Discussion_intro}
% or (if already inside Chapters/Discussion.tex):
%   \input{Discussion_intro}
%
% -----------------------------------------------------------------------------

% === Start writing here ===

% Example scaffold (delete/replace freely):
% \noindent
% This chapter synthesizes the results across RQ1--RQ3 and explains what they
% imply for corpus-independent CLD generation, hallucination detection, and
% evidence-grounded validation. We first restate the core empirical takeaways,
% then interpret failure modes and boundary conditions, and finally discuss
% practical deployment guidance and future work.

\noindent
This chapter synthesizes the results from RQ1--RQ3 on corpus-independent causal loop diagram (CLD) edge generation and hallucination mitigation. We evaluated two LLM-as-a-Judge approaches (correctness-based and citation-based), an LLM-as-a-corrector loop, and generator-side uncertainty quantification (UQ) metrics across three health-domain CLDs. To separate controlled from ecological settings, we test both synthetic ground truth (GT Synth; known corruptions) and literature ground truth (GT Lit; deviation from published expert CLDs). Finally, we explore adaptive test-time compute via an exploratory pilot of Deep Research (DR): multi-agent literature search for direct causal evidence, including a judge-triggered ``smart deployment'' estimate.

\medskip
\noindent
\textbf{Roadmap and framing.}
The discussion that follows focuses less on re-stating each result and more on interpreting \emph{why} performance shifts across settings, and what those shifts imply for how CLD ``hallucinations'' should be operationalised. In particular, we reconcile the strong controlled performance on GT~Synth with the weaker, domain-dependent behaviour on GT~Lit by separating three sources of disagreement: (i) genuine model errors, (ii) retrieval and evidence-type constraints that limit what can be validated from accessible literature, and (iii) construct and scope mismatches between expert CLDs and edge-level evaluation. We then use this lens to analyse failure modes in judge--corrector loops (including Goodhart-like proxy optimisation), to clarify what uncertainty signals can and cannot support in deployment, and to position literature-oriented search (Deep Research) as a form of evidence-gathering test-time compute that may enrich the contested edge space. Throughout, the goal is practical: to outline a workflow where LLMs propose and stress-test candidate edges, automated signals triage uncertainty, and expert judgement is reserved for construct definition and applicability assessment.


% or (if already inside Chapters/Discussion.tex):
%   %==============================================================================
% Discussion Intro (draft scratchpad)
%==============================================================================
%
% Use this file to draft/iterate on the opening of the Discussion chapter.
% It is NOT included anywhere by default.
%
% If you want to include it in the thesis build later, add one of:
%   \input{Chapters/Discussion_intro}
% or (if already inside Chapters/Discussion.tex):
%   \input{Discussion_intro}
%
% -----------------------------------------------------------------------------

% === Start writing here ===

% Example scaffold (delete/replace freely):
% \noindent
% This chapter synthesizes the results across RQ1--RQ3 and explains what they
% imply for corpus-independent CLD generation, hallucination detection, and
% evidence-grounded validation. We first restate the core empirical takeaways,
% then interpret failure modes and boundary conditions, and finally discuss
% practical deployment guidance and future work.

\noindent
This chapter synthesizes the results from RQ1--RQ3 on corpus-independent causal loop diagram (CLD) edge generation and hallucination mitigation. We evaluated two LLM-as-a-Judge approaches (correctness-based and citation-based), an LLM-as-a-corrector loop, and generator-side uncertainty quantification (UQ) metrics across three health-domain CLDs. To separate controlled from ecological settings, we test both synthetic ground truth (GT Synth; known corruptions) and literature ground truth (GT Lit; deviation from published expert CLDs). Finally, we explore adaptive test-time compute via an exploratory pilot of Deep Research (DR): multi-agent literature search for direct causal evidence, including a judge-triggered ``smart deployment'' estimate.

\medskip
\noindent
\textbf{Roadmap and framing.}
The discussion that follows focuses less on re-stating each result and more on interpreting \emph{why} performance shifts across settings, and what those shifts imply for how CLD ``hallucinations'' should be operationalised. In particular, we reconcile the strong controlled performance on GT~Synth with the weaker, domain-dependent behaviour on GT~Lit by separating three sources of disagreement: (i) genuine model errors, (ii) retrieval and evidence-type constraints that limit what can be validated from accessible literature, and (iii) construct and scope mismatches between expert CLDs and edge-level evaluation. We then use this lens to analyse failure modes in judge--corrector loops (including Goodhart-like proxy optimisation), to clarify what uncertainty signals can and cannot support in deployment, and to position literature-oriented search (Deep Research) as a form of evidence-gathering test-time compute that may enrich the contested edge space. Throughout, the goal is practical: to outline a workflow where LLMs propose and stress-test candidate edges, automated signals triage uncertainty, and expert judgement is reserved for construct definition and applicability assessment.


%
% -----------------------------------------------------------------------------

% === Start writing here ===

% Example scaffold (delete/replace freely):
% \noindent
% This chapter synthesizes the results across RQ1--RQ3 and explains what they
% imply for corpus-independent CLD generation, hallucination detection, and
% evidence-grounded validation. We first restate the core empirical takeaways,
% then interpret failure modes and boundary conditions, and finally discuss
% practical deployment guidance and future work.

\noindent
This chapter synthesizes the results from RQ1--RQ3 on corpus-independent causal loop diagram (CLD) edge generation and hallucination mitigation. We evaluated two LLM-as-a-Judge approaches (correctness-based and citation-based), an LLM-as-a-corrector loop, and generator-side uncertainty quantification (UQ) metrics across three health-domain CLDs. To separate controlled from ecological settings, we test both synthetic ground truth (GT Synth; known corruptions) and literature ground truth (GT Lit; deviation from published expert CLDs). Finally, we explore adaptive test-time compute via an exploratory pilot of Deep Research (DR): multi-agent literature search for direct causal evidence, including a judge-triggered ``smart deployment'' estimate.

\medskip
\noindent
\textbf{Roadmap and framing.}
The discussion that follows focuses less on re-stating each result and more on interpreting \emph{why} performance shifts across settings, and what those shifts imply for how CLD ``hallucinations'' should be operationalised. In particular, we reconcile the strong controlled performance on GT~Synth with the weaker, domain-dependent behaviour on GT~Lit by separating three sources of disagreement: (i) genuine model errors, (ii) retrieval and evidence-type constraints that limit what can be validated from accessible literature, and (iii) construct and scope mismatches between expert CLDs and edge-level evaluation. We then use this lens to analyse failure modes in judge--corrector loops (including Goodhart-like proxy optimisation), to clarify what uncertainty signals can and cannot support in deployment, and to position literature-oriented search (Deep Research) as a form of evidence-gathering test-time compute that may enrich the contested edge space. Throughout, the goal is practical: to outline a workflow where LLMs propose and stress-test candidate edges, automated signals triage uncertainty, and expert judgement is reserved for construct definition and applicability assessment.


% or (if already inside Chapters/Discussion.tex):
%   %==============================================================================
% Discussion Intro (draft scratchpad)
%==============================================================================
%
% Use this file to draft/iterate on the opening of the Discussion chapter.
% It is NOT included anywhere by default.
%
% If you want to include it in the thesis build later, add one of:
%   %==============================================================================
% Discussion Intro (draft scratchpad)
%==============================================================================
%
% Use this file to draft/iterate on the opening of the Discussion chapter.
% It is NOT included anywhere by default.
%
% If you want to include it in the thesis build later, add one of:
%   \input{Chapters/Discussion_intro}
% or (if already inside Chapters/Discussion.tex):
%   \input{Discussion_intro}
%
% -----------------------------------------------------------------------------

% === Start writing here ===

% Example scaffold (delete/replace freely):
% \noindent
% This chapter synthesizes the results across RQ1--RQ3 and explains what they
% imply for corpus-independent CLD generation, hallucination detection, and
% evidence-grounded validation. We first restate the core empirical takeaways,
% then interpret failure modes and boundary conditions, and finally discuss
% practical deployment guidance and future work.

\noindent
This chapter synthesizes the results from RQ1--RQ3 on corpus-independent causal loop diagram (CLD) edge generation and hallucination mitigation. We evaluated two LLM-as-a-Judge approaches (correctness-based and citation-based), an LLM-as-a-corrector loop, and generator-side uncertainty quantification (UQ) metrics across three health-domain CLDs. To separate controlled from ecological settings, we test both synthetic ground truth (GT Synth; known corruptions) and literature ground truth (GT Lit; deviation from published expert CLDs). Finally, we explore adaptive test-time compute via an exploratory pilot of Deep Research (DR): multi-agent literature search for direct causal evidence, including a judge-triggered ``smart deployment'' estimate.

\medskip
\noindent
\textbf{Roadmap and framing.}
The discussion that follows focuses less on re-stating each result and more on interpreting \emph{why} performance shifts across settings, and what those shifts imply for how CLD ``hallucinations'' should be operationalised. In particular, we reconcile the strong controlled performance on GT~Synth with the weaker, domain-dependent behaviour on GT~Lit by separating three sources of disagreement: (i) genuine model errors, (ii) retrieval and evidence-type constraints that limit what can be validated from accessible literature, and (iii) construct and scope mismatches between expert CLDs and edge-level evaluation. We then use this lens to analyse failure modes in judge--corrector loops (including Goodhart-like proxy optimisation), to clarify what uncertainty signals can and cannot support in deployment, and to position literature-oriented search (Deep Research) as a form of evidence-gathering test-time compute that may enrich the contested edge space. Throughout, the goal is practical: to outline a workflow where LLMs propose and stress-test candidate edges, automated signals triage uncertainty, and expert judgement is reserved for construct definition and applicability assessment.


% or (if already inside Chapters/Discussion.tex):
%   %==============================================================================
% Discussion Intro (draft scratchpad)
%==============================================================================
%
% Use this file to draft/iterate on the opening of the Discussion chapter.
% It is NOT included anywhere by default.
%
% If you want to include it in the thesis build later, add one of:
%   \input{Chapters/Discussion_intro}
% or (if already inside Chapters/Discussion.tex):
%   \input{Discussion_intro}
%
% -----------------------------------------------------------------------------

% === Start writing here ===

% Example scaffold (delete/replace freely):
% \noindent
% This chapter synthesizes the results across RQ1--RQ3 and explains what they
% imply for corpus-independent CLD generation, hallucination detection, and
% evidence-grounded validation. We first restate the core empirical takeaways,
% then interpret failure modes and boundary conditions, and finally discuss
% practical deployment guidance and future work.

\noindent
This chapter synthesizes the results from RQ1--RQ3 on corpus-independent causal loop diagram (CLD) edge generation and hallucination mitigation. We evaluated two LLM-as-a-Judge approaches (correctness-based and citation-based), an LLM-as-a-corrector loop, and generator-side uncertainty quantification (UQ) metrics across three health-domain CLDs. To separate controlled from ecological settings, we test both synthetic ground truth (GT Synth; known corruptions) and literature ground truth (GT Lit; deviation from published expert CLDs). Finally, we explore adaptive test-time compute via an exploratory pilot of Deep Research (DR): multi-agent literature search for direct causal evidence, including a judge-triggered ``smart deployment'' estimate.

\medskip
\noindent
\textbf{Roadmap and framing.}
The discussion that follows focuses less on re-stating each result and more on interpreting \emph{why} performance shifts across settings, and what those shifts imply for how CLD ``hallucinations'' should be operationalised. In particular, we reconcile the strong controlled performance on GT~Synth with the weaker, domain-dependent behaviour on GT~Lit by separating three sources of disagreement: (i) genuine model errors, (ii) retrieval and evidence-type constraints that limit what can be validated from accessible literature, and (iii) construct and scope mismatches between expert CLDs and edge-level evaluation. We then use this lens to analyse failure modes in judge--corrector loops (including Goodhart-like proxy optimisation), to clarify what uncertainty signals can and cannot support in deployment, and to position literature-oriented search (Deep Research) as a form of evidence-gathering test-time compute that may enrich the contested edge space. Throughout, the goal is practical: to outline a workflow where LLMs propose and stress-test candidate edges, automated signals triage uncertainty, and expert judgement is reserved for construct definition and applicability assessment.


%
% -----------------------------------------------------------------------------

% === Start writing here ===

% Example scaffold (delete/replace freely):
% \noindent
% This chapter synthesizes the results across RQ1--RQ3 and explains what they
% imply for corpus-independent CLD generation, hallucination detection, and
% evidence-grounded validation. We first restate the core empirical takeaways,
% then interpret failure modes and boundary conditions, and finally discuss
% practical deployment guidance and future work.

\noindent
This chapter synthesizes the results from RQ1--RQ3 on corpus-independent causal loop diagram (CLD) edge generation and hallucination mitigation. We evaluated two LLM-as-a-Judge approaches (correctness-based and citation-based), an LLM-as-a-corrector loop, and generator-side uncertainty quantification (UQ) metrics across three health-domain CLDs. To separate controlled from ecological settings, we test both synthetic ground truth (GT Synth; known corruptions) and literature ground truth (GT Lit; deviation from published expert CLDs). Finally, we explore adaptive test-time compute via an exploratory pilot of Deep Research (DR): multi-agent literature search for direct causal evidence, including a judge-triggered ``smart deployment'' estimate.

\medskip
\noindent
\textbf{Roadmap and framing.}
The discussion that follows focuses less on re-stating each result and more on interpreting \emph{why} performance shifts across settings, and what those shifts imply for how CLD ``hallucinations'' should be operationalised. In particular, we reconcile the strong controlled performance on GT~Synth with the weaker, domain-dependent behaviour on GT~Lit by separating three sources of disagreement: (i) genuine model errors, (ii) retrieval and evidence-type constraints that limit what can be validated from accessible literature, and (iii) construct and scope mismatches between expert CLDs and edge-level evaluation. We then use this lens to analyse failure modes in judge--corrector loops (including Goodhart-like proxy optimisation), to clarify what uncertainty signals can and cannot support in deployment, and to position literature-oriented search (Deep Research) as a form of evidence-gathering test-time compute that may enrich the contested edge space. Throughout, the goal is practical: to outline a workflow where LLMs propose and stress-test candidate edges, automated signals triage uncertainty, and expert judgement is reserved for construct definition and applicability assessment.


%
% -----------------------------------------------------------------------------

% === Start writing here ===

% Example scaffold (delete/replace freely):
% \noindent
% This chapter synthesizes the results across RQ1--RQ3 and explains what they
% imply for corpus-independent CLD generation, hallucination detection, and
% evidence-grounded validation. We first restate the core empirical takeaways,
% then interpret failure modes and boundary conditions, and finally discuss
% practical deployment guidance and future work.

\noindent
This chapter synthesizes the results from RQ1--RQ3 on corpus-independent causal loop diagram (CLD) edge generation and hallucination mitigation. We evaluated two LLM-as-a-Judge approaches (correctness-based and citation-based), an LLM-as-a-corrector loop, and generator-side uncertainty quantification (UQ) metrics across three health-domain CLDs. To separate controlled from ecological settings, we test both synthetic ground truth (GT Synth; known corruptions) and literature ground truth (GT Lit; deviation from published expert CLDs). Finally, we explore adaptive test-time compute via an exploratory pilot of Deep Research (DR): multi-agent literature search for direct causal evidence, including a judge-triggered ``smart deployment'' estimate.

\medskip
\noindent
\textbf{Roadmap and framing.}
The discussion that follows focuses less on re-stating each result and more on interpreting \emph{why} performance shifts across settings, and what those shifts imply for how CLD ``hallucinations'' should be operationalised. In particular, we reconcile the strong controlled performance on GT~Synth with the weaker, domain-dependent behaviour on GT~Lit by separating three sources of disagreement: (i) genuine model errors, (ii) retrieval and evidence-type constraints that limit what can be validated from accessible literature, and (iii) construct and scope mismatches between expert CLDs and edge-level evaluation. We then use this lens to analyse failure modes in judge--corrector loops (including Goodhart-like proxy optimisation), to clarify what uncertainty signals can and cannot support in deployment, and to position literature-oriented search (Deep Research) as a form of evidence-gathering test-time compute that may enrich the contested edge space. Throughout, the goal is practical: to outline a workflow where LLMs propose and stress-test candidate edges, automated signals triage uncertainty, and expert judgement is reserved for construct definition and applicability assessment.


% or (if already inside Chapters/Discussion.tex):
%   %==============================================================================
% Discussion Intro (draft scratchpad)
%==============================================================================
%
% Use this file to draft/iterate on the opening of the Discussion chapter.
% It is NOT included anywhere by default.
%
% If you want to include it in the thesis build later, add one of:
%   %==============================================================================
% Discussion Intro (draft scratchpad)
%==============================================================================
%
% Use this file to draft/iterate on the opening of the Discussion chapter.
% It is NOT included anywhere by default.
%
% If you want to include it in the thesis build later, add one of:
%   %==============================================================================
% Discussion Intro (draft scratchpad)
%==============================================================================
%
% Use this file to draft/iterate on the opening of the Discussion chapter.
% It is NOT included anywhere by default.
%
% If you want to include it in the thesis build later, add one of:
%   \input{Chapters/Discussion_intro}
% or (if already inside Chapters/Discussion.tex):
%   \input{Discussion_intro}
%
% -----------------------------------------------------------------------------

% === Start writing here ===

% Example scaffold (delete/replace freely):
% \noindent
% This chapter synthesizes the results across RQ1--RQ3 and explains what they
% imply for corpus-independent CLD generation, hallucination detection, and
% evidence-grounded validation. We first restate the core empirical takeaways,
% then interpret failure modes and boundary conditions, and finally discuss
% practical deployment guidance and future work.

\noindent
This chapter synthesizes the results from RQ1--RQ3 on corpus-independent causal loop diagram (CLD) edge generation and hallucination mitigation. We evaluated two LLM-as-a-Judge approaches (correctness-based and citation-based), an LLM-as-a-corrector loop, and generator-side uncertainty quantification (UQ) metrics across three health-domain CLDs. To separate controlled from ecological settings, we test both synthetic ground truth (GT Synth; known corruptions) and literature ground truth (GT Lit; deviation from published expert CLDs). Finally, we explore adaptive test-time compute via an exploratory pilot of Deep Research (DR): multi-agent literature search for direct causal evidence, including a judge-triggered ``smart deployment'' estimate.

\medskip
\noindent
\textbf{Roadmap and framing.}
The discussion that follows focuses less on re-stating each result and more on interpreting \emph{why} performance shifts across settings, and what those shifts imply for how CLD ``hallucinations'' should be operationalised. In particular, we reconcile the strong controlled performance on GT~Synth with the weaker, domain-dependent behaviour on GT~Lit by separating three sources of disagreement: (i) genuine model errors, (ii) retrieval and evidence-type constraints that limit what can be validated from accessible literature, and (iii) construct and scope mismatches between expert CLDs and edge-level evaluation. We then use this lens to analyse failure modes in judge--corrector loops (including Goodhart-like proxy optimisation), to clarify what uncertainty signals can and cannot support in deployment, and to position literature-oriented search (Deep Research) as a form of evidence-gathering test-time compute that may enrich the contested edge space. Throughout, the goal is practical: to outline a workflow where LLMs propose and stress-test candidate edges, automated signals triage uncertainty, and expert judgement is reserved for construct definition and applicability assessment.


% or (if already inside Chapters/Discussion.tex):
%   %==============================================================================
% Discussion Intro (draft scratchpad)
%==============================================================================
%
% Use this file to draft/iterate on the opening of the Discussion chapter.
% It is NOT included anywhere by default.
%
% If you want to include it in the thesis build later, add one of:
%   \input{Chapters/Discussion_intro}
% or (if already inside Chapters/Discussion.tex):
%   \input{Discussion_intro}
%
% -----------------------------------------------------------------------------

% === Start writing here ===

% Example scaffold (delete/replace freely):
% \noindent
% This chapter synthesizes the results across RQ1--RQ3 and explains what they
% imply for corpus-independent CLD generation, hallucination detection, and
% evidence-grounded validation. We first restate the core empirical takeaways,
% then interpret failure modes and boundary conditions, and finally discuss
% practical deployment guidance and future work.

\noindent
This chapter synthesizes the results from RQ1--RQ3 on corpus-independent causal loop diagram (CLD) edge generation and hallucination mitigation. We evaluated two LLM-as-a-Judge approaches (correctness-based and citation-based), an LLM-as-a-corrector loop, and generator-side uncertainty quantification (UQ) metrics across three health-domain CLDs. To separate controlled from ecological settings, we test both synthetic ground truth (GT Synth; known corruptions) and literature ground truth (GT Lit; deviation from published expert CLDs). Finally, we explore adaptive test-time compute via an exploratory pilot of Deep Research (DR): multi-agent literature search for direct causal evidence, including a judge-triggered ``smart deployment'' estimate.

\medskip
\noindent
\textbf{Roadmap and framing.}
The discussion that follows focuses less on re-stating each result and more on interpreting \emph{why} performance shifts across settings, and what those shifts imply for how CLD ``hallucinations'' should be operationalised. In particular, we reconcile the strong controlled performance on GT~Synth with the weaker, domain-dependent behaviour on GT~Lit by separating three sources of disagreement: (i) genuine model errors, (ii) retrieval and evidence-type constraints that limit what can be validated from accessible literature, and (iii) construct and scope mismatches between expert CLDs and edge-level evaluation. We then use this lens to analyse failure modes in judge--corrector loops (including Goodhart-like proxy optimisation), to clarify what uncertainty signals can and cannot support in deployment, and to position literature-oriented search (Deep Research) as a form of evidence-gathering test-time compute that may enrich the contested edge space. Throughout, the goal is practical: to outline a workflow where LLMs propose and stress-test candidate edges, automated signals triage uncertainty, and expert judgement is reserved for construct definition and applicability assessment.


%
% -----------------------------------------------------------------------------

% === Start writing here ===

% Example scaffold (delete/replace freely):
% \noindent
% This chapter synthesizes the results across RQ1--RQ3 and explains what they
% imply for corpus-independent CLD generation, hallucination detection, and
% evidence-grounded validation. We first restate the core empirical takeaways,
% then interpret failure modes and boundary conditions, and finally discuss
% practical deployment guidance and future work.

\noindent
This chapter synthesizes the results from RQ1--RQ3 on corpus-independent causal loop diagram (CLD) edge generation and hallucination mitigation. We evaluated two LLM-as-a-Judge approaches (correctness-based and citation-based), an LLM-as-a-corrector loop, and generator-side uncertainty quantification (UQ) metrics across three health-domain CLDs. To separate controlled from ecological settings, we test both synthetic ground truth (GT Synth; known corruptions) and literature ground truth (GT Lit; deviation from published expert CLDs). Finally, we explore adaptive test-time compute via an exploratory pilot of Deep Research (DR): multi-agent literature search for direct causal evidence, including a judge-triggered ``smart deployment'' estimate.

\medskip
\noindent
\textbf{Roadmap and framing.}
The discussion that follows focuses less on re-stating each result and more on interpreting \emph{why} performance shifts across settings, and what those shifts imply for how CLD ``hallucinations'' should be operationalised. In particular, we reconcile the strong controlled performance on GT~Synth with the weaker, domain-dependent behaviour on GT~Lit by separating three sources of disagreement: (i) genuine model errors, (ii) retrieval and evidence-type constraints that limit what can be validated from accessible literature, and (iii) construct and scope mismatches between expert CLDs and edge-level evaluation. We then use this lens to analyse failure modes in judge--corrector loops (including Goodhart-like proxy optimisation), to clarify what uncertainty signals can and cannot support in deployment, and to position literature-oriented search (Deep Research) as a form of evidence-gathering test-time compute that may enrich the contested edge space. Throughout, the goal is practical: to outline a workflow where LLMs propose and stress-test candidate edges, automated signals triage uncertainty, and expert judgement is reserved for construct definition and applicability assessment.


% or (if already inside Chapters/Discussion.tex):
%   %==============================================================================
% Discussion Intro (draft scratchpad)
%==============================================================================
%
% Use this file to draft/iterate on the opening of the Discussion chapter.
% It is NOT included anywhere by default.
%
% If you want to include it in the thesis build later, add one of:
%   %==============================================================================
% Discussion Intro (draft scratchpad)
%==============================================================================
%
% Use this file to draft/iterate on the opening of the Discussion chapter.
% It is NOT included anywhere by default.
%
% If you want to include it in the thesis build later, add one of:
%   \input{Chapters/Discussion_intro}
% or (if already inside Chapters/Discussion.tex):
%   \input{Discussion_intro}
%
% -----------------------------------------------------------------------------

% === Start writing here ===

% Example scaffold (delete/replace freely):
% \noindent
% This chapter synthesizes the results across RQ1--RQ3 and explains what they
% imply for corpus-independent CLD generation, hallucination detection, and
% evidence-grounded validation. We first restate the core empirical takeaways,
% then interpret failure modes and boundary conditions, and finally discuss
% practical deployment guidance and future work.

\noindent
This chapter synthesizes the results from RQ1--RQ3 on corpus-independent causal loop diagram (CLD) edge generation and hallucination mitigation. We evaluated two LLM-as-a-Judge approaches (correctness-based and citation-based), an LLM-as-a-corrector loop, and generator-side uncertainty quantification (UQ) metrics across three health-domain CLDs. To separate controlled from ecological settings, we test both synthetic ground truth (GT Synth; known corruptions) and literature ground truth (GT Lit; deviation from published expert CLDs). Finally, we explore adaptive test-time compute via an exploratory pilot of Deep Research (DR): multi-agent literature search for direct causal evidence, including a judge-triggered ``smart deployment'' estimate.

\medskip
\noindent
\textbf{Roadmap and framing.}
The discussion that follows focuses less on re-stating each result and more on interpreting \emph{why} performance shifts across settings, and what those shifts imply for how CLD ``hallucinations'' should be operationalised. In particular, we reconcile the strong controlled performance on GT~Synth with the weaker, domain-dependent behaviour on GT~Lit by separating three sources of disagreement: (i) genuine model errors, (ii) retrieval and evidence-type constraints that limit what can be validated from accessible literature, and (iii) construct and scope mismatches between expert CLDs and edge-level evaluation. We then use this lens to analyse failure modes in judge--corrector loops (including Goodhart-like proxy optimisation), to clarify what uncertainty signals can and cannot support in deployment, and to position literature-oriented search (Deep Research) as a form of evidence-gathering test-time compute that may enrich the contested edge space. Throughout, the goal is practical: to outline a workflow where LLMs propose and stress-test candidate edges, automated signals triage uncertainty, and expert judgement is reserved for construct definition and applicability assessment.


% or (if already inside Chapters/Discussion.tex):
%   %==============================================================================
% Discussion Intro (draft scratchpad)
%==============================================================================
%
% Use this file to draft/iterate on the opening of the Discussion chapter.
% It is NOT included anywhere by default.
%
% If you want to include it in the thesis build later, add one of:
%   \input{Chapters/Discussion_intro}
% or (if already inside Chapters/Discussion.tex):
%   \input{Discussion_intro}
%
% -----------------------------------------------------------------------------

% === Start writing here ===

% Example scaffold (delete/replace freely):
% \noindent
% This chapter synthesizes the results across RQ1--RQ3 and explains what they
% imply for corpus-independent CLD generation, hallucination detection, and
% evidence-grounded validation. We first restate the core empirical takeaways,
% then interpret failure modes and boundary conditions, and finally discuss
% practical deployment guidance and future work.

\noindent
This chapter synthesizes the results from RQ1--RQ3 on corpus-independent causal loop diagram (CLD) edge generation and hallucination mitigation. We evaluated two LLM-as-a-Judge approaches (correctness-based and citation-based), an LLM-as-a-corrector loop, and generator-side uncertainty quantification (UQ) metrics across three health-domain CLDs. To separate controlled from ecological settings, we test both synthetic ground truth (GT Synth; known corruptions) and literature ground truth (GT Lit; deviation from published expert CLDs). Finally, we explore adaptive test-time compute via an exploratory pilot of Deep Research (DR): multi-agent literature search for direct causal evidence, including a judge-triggered ``smart deployment'' estimate.

\medskip
\noindent
\textbf{Roadmap and framing.}
The discussion that follows focuses less on re-stating each result and more on interpreting \emph{why} performance shifts across settings, and what those shifts imply for how CLD ``hallucinations'' should be operationalised. In particular, we reconcile the strong controlled performance on GT~Synth with the weaker, domain-dependent behaviour on GT~Lit by separating three sources of disagreement: (i) genuine model errors, (ii) retrieval and evidence-type constraints that limit what can be validated from accessible literature, and (iii) construct and scope mismatches between expert CLDs and edge-level evaluation. We then use this lens to analyse failure modes in judge--corrector loops (including Goodhart-like proxy optimisation), to clarify what uncertainty signals can and cannot support in deployment, and to position literature-oriented search (Deep Research) as a form of evidence-gathering test-time compute that may enrich the contested edge space. Throughout, the goal is practical: to outline a workflow where LLMs propose and stress-test candidate edges, automated signals triage uncertainty, and expert judgement is reserved for construct definition and applicability assessment.


%
% -----------------------------------------------------------------------------

% === Start writing here ===

% Example scaffold (delete/replace freely):
% \noindent
% This chapter synthesizes the results across RQ1--RQ3 and explains what they
% imply for corpus-independent CLD generation, hallucination detection, and
% evidence-grounded validation. We first restate the core empirical takeaways,
% then interpret failure modes and boundary conditions, and finally discuss
% practical deployment guidance and future work.

\noindent
This chapter synthesizes the results from RQ1--RQ3 on corpus-independent causal loop diagram (CLD) edge generation and hallucination mitigation. We evaluated two LLM-as-a-Judge approaches (correctness-based and citation-based), an LLM-as-a-corrector loop, and generator-side uncertainty quantification (UQ) metrics across three health-domain CLDs. To separate controlled from ecological settings, we test both synthetic ground truth (GT Synth; known corruptions) and literature ground truth (GT Lit; deviation from published expert CLDs). Finally, we explore adaptive test-time compute via an exploratory pilot of Deep Research (DR): multi-agent literature search for direct causal evidence, including a judge-triggered ``smart deployment'' estimate.

\medskip
\noindent
\textbf{Roadmap and framing.}
The discussion that follows focuses less on re-stating each result and more on interpreting \emph{why} performance shifts across settings, and what those shifts imply for how CLD ``hallucinations'' should be operationalised. In particular, we reconcile the strong controlled performance on GT~Synth with the weaker, domain-dependent behaviour on GT~Lit by separating three sources of disagreement: (i) genuine model errors, (ii) retrieval and evidence-type constraints that limit what can be validated from accessible literature, and (iii) construct and scope mismatches between expert CLDs and edge-level evaluation. We then use this lens to analyse failure modes in judge--corrector loops (including Goodhart-like proxy optimisation), to clarify what uncertainty signals can and cannot support in deployment, and to position literature-oriented search (Deep Research) as a form of evidence-gathering test-time compute that may enrich the contested edge space. Throughout, the goal is practical: to outline a workflow where LLMs propose and stress-test candidate edges, automated signals triage uncertainty, and expert judgement is reserved for construct definition and applicability assessment.


%
% -----------------------------------------------------------------------------

% === Start writing here ===

% Example scaffold (delete/replace freely):
% \noindent
% This chapter synthesizes the results across RQ1--RQ3 and explains what they
% imply for corpus-independent CLD generation, hallucination detection, and
% evidence-grounded validation. We first restate the core empirical takeaways,
% then interpret failure modes and boundary conditions, and finally discuss
% practical deployment guidance and future work.

\noindent
This chapter synthesizes the results from RQ1--RQ3 on corpus-independent causal loop diagram (CLD) edge generation and hallucination mitigation. We evaluated two LLM-as-a-Judge approaches (correctness-based and citation-based), an LLM-as-a-corrector loop, and generator-side uncertainty quantification (UQ) metrics across three health-domain CLDs. To separate controlled from ecological settings, we test both synthetic ground truth (GT Synth; known corruptions) and literature ground truth (GT Lit; deviation from published expert CLDs). Finally, we explore adaptive test-time compute via an exploratory pilot of Deep Research (DR): multi-agent literature search for direct causal evidence, including a judge-triggered ``smart deployment'' estimate.

\medskip
\noindent
\textbf{Roadmap and framing.}
The discussion that follows focuses less on re-stating each result and more on interpreting \emph{why} performance shifts across settings, and what those shifts imply for how CLD ``hallucinations'' should be operationalised. In particular, we reconcile the strong controlled performance on GT~Synth with the weaker, domain-dependent behaviour on GT~Lit by separating three sources of disagreement: (i) genuine model errors, (ii) retrieval and evidence-type constraints that limit what can be validated from accessible literature, and (iii) construct and scope mismatches between expert CLDs and edge-level evaluation. We then use this lens to analyse failure modes in judge--corrector loops (including Goodhart-like proxy optimisation), to clarify what uncertainty signals can and cannot support in deployment, and to position literature-oriented search (Deep Research) as a form of evidence-gathering test-time compute that may enrich the contested edge space. Throughout, the goal is practical: to outline a workflow where LLMs propose and stress-test candidate edges, automated signals triage uncertainty, and expert judgement is reserved for construct definition and applicability assessment.


%
% -----------------------------------------------------------------------------

% === Start writing here ===

% Example scaffold (delete/replace freely):
% \noindent
% This chapter synthesizes the results across RQ1--RQ3 and explains what they
% imply for corpus-independent CLD generation, hallucination detection, and
% evidence-grounded validation. We first restate the core empirical takeaways,
% then interpret failure modes and boundary conditions, and finally discuss
% practical deployment guidance and future work.

\noindent
This chapter synthesizes the results from RQ1--RQ3 on corpus-independent causal loop diagram (CLD) edge generation and hallucination mitigation. We evaluated two LLM-as-a-Judge approaches (correctness-based and citation-based), an LLM-as-a-corrector loop, and generator-side uncertainty quantification (UQ) metrics across three health-domain CLDs. To separate controlled from ecological settings, we test both synthetic ground truth (GT Synth; known corruptions) and literature ground truth (GT Lit; deviation from published expert CLDs). Finally, we explore adaptive test-time compute via an exploratory pilot of Deep Research (DR): multi-agent literature search for direct causal evidence, including a judge-triggered ``smart deployment'' estimate.

\medskip
\noindent
\textbf{Roadmap and framing.}
The discussion that follows focuses less on re-stating each result and more on interpreting \emph{why} performance shifts across settings, and what those shifts imply for how CLD ``hallucinations'' should be operationalised. In particular, we reconcile the strong controlled performance on GT~Synth with the weaker, domain-dependent behaviour on GT~Lit by separating three sources of disagreement: (i) genuine model errors, (ii) retrieval and evidence-type constraints that limit what can be validated from accessible literature, and (iii) construct and scope mismatches between expert CLDs and edge-level evaluation. We then use this lens to analyse failure modes in judge--corrector loops (including Goodhart-like proxy optimisation), to clarify what uncertainty signals can and cannot support in deployment, and to position literature-oriented search (Deep Research) as a form of evidence-gathering test-time compute that may enrich the contested edge space. Throughout, the goal is practical: to outline a workflow where LLMs propose and stress-test candidate edges, automated signals triage uncertainty, and expert judgement is reserved for construct definition and applicability assessment.

